\documentclass[[a4paper,11pt]{article}
\renewcommand{\rmdefault}{ptm}
\usepackage[scaled=0.92]{helvet}
\usepackage{courier,xcolor,colortbl,listings,parskip,graphicx,fancyvrb,fancyhdr,lastpage}
\usepackage{float,framed}
\normalfont
\usepackage[T1]{fontenc}
\setlength{\parskip}{7pt}
\usepackage[toc,page]{appendix}
\usepackage[hmargin=2.5cm,vmargin=2cm]{geometry}
\usepackage[utf8]{inputenc}
\usepackage[brazil]{babel}
\pagestyle{fancy}
\setlength{\headheight}{120pt}
\setlength{\headsep}{30pt}
\setlength{\textheight}{550pt}
\renewcommand{\headrulewidth}{0pt}
\lhead{}
\rhead{}
\chead{\includegraphics{brasao.jpg}\\
        \large \textbf{PRESIDÊNCIA DA REPÚBLICA}\\
        \large SECRETARIA-GERAL\\
        \large Secretaria-Executiva}
\cfoot{}
\rfoot{\thepage /\pageref{LastPage}}
\hyphenation{par-ti-ci-pa-ção}
\bibliographystyle{ieeetr}

\newcommand{\MyName}{Daniela Soares Feitosa}
\newcommand{\MyEmail}{daniela@colivre.coop.br}
\newcommand{\ContractNumber}{2013/000292}
\newcommand{\ProjectCode}{Projeto PNUD BRA/12/018}
\newcommand{\NomeSecretaria}{Secretaria Geral da Presidência da República}
\newcommand{\SiglaSecretaria}{SG/PR}
\newcommand{\ProductNumber}{05}
\newcommand{\ProductDescription}{Documento com proposta para 
desenvolvimento do código do tema das
comunidades contendo exemplos e códigos.}
\newcommand{\MesEntrega}{Maio de 2014}
\newcommand{\DiaEntrega}{26}

\begin{document}
\lstset{language=Ruby}
\definecolor{light-gray}{gray}{0.95}
\lstdefinestyle{codeFrame}{backgroundcolor=\color{light-gray},frame=lines}

\textbf{\ProjectCode \ -} \ProductDescription

\vspace{3cm}

\begin{minipage}{0.5\textwidth}
  \textbf{Consultora: \MyName}
  \newline
  \textbf{Contrato nº: \ContractNumber}
  \newline
  \textbf{Produto / nº: \ProductNumber}
\end{minipage}

\vspace{2cm}

\textbf{Assinatura do Consultor}

\begin{framed}
Local e data: Brasília/DF, \line(1,0){20} \ de \line(1,0){100} \ de 2014
\newline
\newline
Assinatura do Consultor: \line(1,0){300}
\end{framed}

\vspace{1cm}

\textbf{Assinatura do Supervisor}

\begin{framed}
Atesto que os serviços foram prestados conforme estabelecido no Contrato
de Consultoria.
\newline
\newline
Local e data: Brasília/DF, \line(1,0){20} \ de \line(1,0){100} \ de 2014
\newline
\newline
Assinatura e Carimbo: \line(1,0){300}
\end{framed}

\clearpage
\newcolumntype{g}{>{\columncolor{light-gray}}l}

\begin{center}
  \begin{tabular}{| g | p{10cm} |}
    \hline
    \textbf{Título} & \ProductDescription \\ \hline
    \textbf{Língua do documento} & Português - Brasil \\ \hline
    \textbf{Documentação de referência} & Português \\ \hline
    \textbf{Unidade responsável} & \NomeSecretaria \
(\SiglaSecretaria) \\ \hline
    \textbf{Criador} & \MyName - \MyEmail \\ \hline
    \textbf{Taxonomias} & Desenvolvimento \\ \hline
    \textbf{Data de aprovação} &  \\ \hline
    \textbf{Público} & \SiglaSecretaria, Parceiros e Sociedade
Civil \\ \hline
    \textbf{Faz parte do} & \ProjectCode \\ \hline
    \textbf{Em conformidade com a} & \NomeSecretaria \\ \hline
    \textbf{Documentos anexos} & Página da Comunidade OSC - Organizações da
Sociedade Civil; CSS do cabeçalho do tema de comunidade; Código javascript para
incluir uma classe identificadora; Código CSS para estilizar os blocos com o padrão
colorido \\ \hline
    \textbf{Revisado em} &  \\ \hline
  \end{tabular}
\end{center}

\clearpage

\tableofcontents
\clearpage
\listoffigures

\clearpage

\section{Introdução}

A participação social no Brasil representa princípio
jurídico-institucional
presente na Constituição Federal de 1988, que a definiu como forma de
afirmação da
democracia e da consolidação da cidadania. Ao incorporar esse princípio
como
referência para a gestão pública, o Governo Federal aprimora os
processos de interação
do Estado com a sociedade e cria as condições institucionais para a
prática da
democracia participativa. Com isso, verifica-se que, além da crescente
participação
social nas decisões governamentais, as políticas públicas ganham maior
legitimidade,
uma vez que expressam as atuais condições socioeconômicas e culturais da
população
brasileira em suas diferentes realidades regionais.

Na estrutura administrativa do Poder Executivo Federal, cabe à
Secretaria-Geral
da Presidência da República (SG/PR) a função de intermediar as relações
do Governo
com as entidades da sociedade civil, conforme competências definidas
pela Lei
10.683/2003 e pelo Decreto no. 7.688/2012. Assim, a SG/PR é órgão
incumbido de
assessorar diretamente a Presidenta da República e os órgãos e entidades
do Governo
Federal no relacionamento e na articulação com os movimentos sociais, o
que inclui a
criação e a implementação de canais que assegurem a consulta e a
participação popular
na discussão e na definição da agenda prioritária do país.

O Brasil tem um rico histórico de efetivação da democracia
participativa, sendo
reconhecido mundialmente. Os instrumentos institucionalizados como
conselhos de
políticas públicas e conferências nacionais foram profundamente
ampliados na última
década, contando com um legado volumoso de práticas e realizações. A
maioria dos
programas de governo já conta com participação social prevista em pelo
menos uma de
suas etapas. As práticas trazidas pelas novas mídias e pela cultura
digital podem
interagir nesses espaços fortalecendo, ampliando e aprofundando a
democracia
participativa, já na abordagem do novo século, pós redes sociais
digitais.

Incentivando os atores a conectar perfis, blogs e demais instâncias de
produção
de conteúdo na rede, o Portal de Participação Social poderá se
estabelecer como um
repositório agregador do conhecimento sobre participação social hoje
disperso na
internet e nas instâncias governamentais. A partir de uma interface
clara que possibilite
a navegação pelos temas, o Portal de Participação Social pode ser um
espaço onde os
diferentes agentes de governo, movimentos, organizações e cidadãos em
geral
encontrarão solo fértil e facilitado para o exercício da pesquisa e do
diálogo.

Tendo por base a premissa de que a incursão e abertura de canais de
acesso ao
poder público na rede aumenta o conhecimento das ações do governo e
diminui as
barreiras para participação de cidadãos, entidades e movimentos, esse
projeto visa a
construção de um conjunto de ferramentas que poderão ser utilizadas por
gestores e
servidores para proporcionar novas formas de participação a serem
apropriadas pela
cidadania. Além disso, esse portal também buscará dar evidência às
formas de
participação existentes no sentido de contextualizar, organizar e
facilitar o acesso do
cidadão às formas de incidir nas diversas etapas das políticas públicas
do governo
brasileiro.

\section{Apresentação}

Este documento é parte da consultoria para o projeto
\textbf{Desenvolvimento de Metodologias de Articulação e Gestão de
Políticas Públicas para Promoção da Democracia Participativa}
(BRA/12/018), firmado entre a Secretaria-Geral da Presidência da República
(SG/PR) e o Programa das Nações Unidas para o Desenvolvimento (PNUD).
A consultoria tem como objetivo a especificação da construção dos códigos das
metodologias de organização da informação e interação participativa do
portal da participação social.

O presente documento apresenta a proposta para
desenvolvimento do código do tema das comunidades contendo exemplos e códigos.
Essa proposta está configurada como produto 5 da consultoria técnica.

\section{Proposta para desenvolvimento do código do tema das 
comunidades contendo exemplos e códigos.}

\newpage

\section{Considerações finais}


\vspace{1cm}

Sem mais nada a acrescentar, coloco-me à disposição.

\vspace{1cm}

\begin{minipage}{\textwidth}
  Brasília/DF, \DiaEntrega \ de \MesEntrega\\[1cm]
  \textbf{\MyName}\\
  \small Consultora do PNUD
\end{minipage}

\newpage
\appendix
\appendixpage
\section{Plugin class} \label{App:PluginCode}





\end{document}
